\chapter{La página web: donde todo se junta}
\label{cap:s-pagina}
La página es la cara del proyecto: es donde la campaña llega al viajero y donde el jurado puede ver los modelos
funcionando. Está publicada en \codigo{catoxp.github.io/torre-del-caribe-unrc-2026}.

\section{Cómo está ordenada}
Brandon pidió que la página empezara por el viaje y dejara los datos al final (``primero hasta arriba el plan de viaje\dots
los datos al final''). Quedó así:

\begin{center}
\small
\begin{tabularx}{\linewidth}{YY}
\encabezado \h{Para el viajero (arriba)} & \h{Los datos (al final)} \\
1. Portada con el planeador: se elige lugar y mes, y la foto cambia con el lugar & 6. ``Los datos'': el problema, el
Radar, cómo llega la gente \\
2. ``Así va a estar'': temporada, gente esperada, clima, riesgo de tormenta y una recomendación & 7. El Pronóstico, los
escenarios y el reparto del dinero \\
3. Qué hacer de día, de tarde y de noche; dónde comer y dónde dormir & 8. La Torre semana a semana y la campaña \\
4. Los lugares, en un mapa 3D & 9. Las 12 fases, las fuentes y quiénes somos \\
5. Experiencias, rutas de 2 y 3 días y preguntas frecuentes & \\
\bottomrule
\end{tabularx}
\normalsize
\end{center}

\section{El planeador: la redistribución que se ve}
\begin{ejemplo}[ · la Ruta en enero de 2027]
Quien elige la Ruta de las pirámides en enero de 2027 ve que es \textbf{temporada alta}: llega 61\,\% más gente que en un mes
promedio y hay 30\,\% de riesgo de pasar del mes más lleno de su historia. Se esperan unos 8,993 visitantes (entre 5,716 y
14,149). La página le recomienda \textbf{Chetumal} ese mes (tranquilo) o \textbf{la Ruta en noviembre}.
\end{ejemplo}

\textbf{Cuándo es temporada alta}: si el mes tiene 20\,\% más gente que un mes normal (forma del año de 1.20 o más) o si hay
10\,\% o más de riesgo de pasar la capacidad probada. Para Cancún y la Riviera Maya, que cambian poco en el año, se usa el
corte del Radar (71.2\,\% de ocupación): así Cancún queda en temporada alta de noviembre a abril.

\textbf{Qué mes recomienda}: el más tranquilo del mismo lugar que no sea temporada alta, que esté fuera de la temporada de
tormentas y que llueva menos que un mes típico.

\begin{error}
\textbf{El mes recomendado era el más lluvioso.} Sin la condición de lluvia, el planeador casi siempre recomendaba
\textbf{junio}: es tranquilo, pero es el mes más lluvioso en Chetumal (195 mm). Sería un mal consejo. Se agregó la condición,
y ahora recomienda noviembre para la Ruta, mayo para la Bahía y febrero para Chetumal.
\end{error}

\textbf{Cancún y la Riviera Maya} aparecen bajo ``¿Ibas al norte?'' con la etiqueta ``Referencia: la campaña no lo
promueve''. Si su mes está lleno, la página \textbf{siempre} recomienda un lugar del sur, con un aviso que no se pierde al
bajar. Nunca recomienda ir al norte.

\section{Qué hacer, dónde comer y dónde dormir}
Los negocios salen del DENUE (oficiales, con nombre y coordenadas), ordenados por cercanía y clasificados por su nombre
(capítulo~\ref{cap:s-patrones}). Cada tarjeta tiene dos botones que abren Google Maps en otra pestaña: ``Cómo llegar'' y
``Reseñas''. La página \textbf{no copia ni descarga nada de Google}: es un enlace público, como compartir una dirección.

\begin{porque}
Brandon había pedido ``recomendarlo mediante Google Maps y webscraping''. El \emph{scraping} (extraer datos de una página con
un programa) de Google Maps lo prohíben sus términos de uso y la regla 4 del proyecto. La API oficial de Google pide una
tarjeta de crédito, deja de funcionar sin internet y su llave quedaría expuesta en un repositorio público. Con el DENUE y un
enlace, la página sigue funcionando sin internet hasta que alguien da clic.
\end{porque}

\section{Las fotos: todas comprobadas}
\begin{error}
\textbf{Las fotos de comida no eran de Quintana Roo.} Se había hecho una galería de 24 platillos con fotos de Wikimedia
Commons. Solo se había revisado que no mencionaran lugares excluidos, pero eran de Mérida, Campeche y otros lugares. Brandon
lo detectó (``que sean fotos de las zonas que sí sean de Quintana Roo''). La galería se retiró completa, y desde entonces
\textbf{toda foto debe tener coordenada GPS dentro del municipio de su lugar}; si no la tiene, el programa se detiene. Hoy
hay 28 fotos de los cinco lugares, cada una con su crédito y su licencia.
\end{error}

\section{Día y noche, y 9 idiomas}
Un botón cambia la página a modo noche (con una animación de sol que se vuelve luna) y otro cambia el idioma: español,
inglés, francés, alemán, italiano, portugués, japonés, coreano y chino. El maya yucateco está preparado, pero no se publica
hasta que lo revise una persona hablante. La parte de ``Los datos'' queda en español.

\section{Una página que calcula}
El plan pedía una web real con un sistema que \textbf{calcula}, no una página con números fijos. Hay dos formas de verla:
\begin{itemize}
\item \textbf{En internet} (GitHub Pages): la página muestra los resultados ya calculados y no hace ninguna llamada a otros
  sitios.
\item \textbf{En la computadora del proyecto}, con el servidor encendido, aparecen dos controles más: ``Pruébalo tú'',
  donde se mueve el presupuesto, la conversión o los topes, y el servidor \textbf{resuelve el modelo de optimización en
  vivo} (tardó 274 milésimas de segundo); y ``Escuchar en vivo'', donde la Torre avanza semana a semana. El servidor se
  enciende con el comando \texttt{python -m torre.api.servidor}.
\end{itemize}
El servidor (FastAPI) tiene 8 direcciones de consulta. La consulta a la base de datos es \textbf{de solo lectura}: rechaza
cualquier intento de cambiar datos o leer archivos (se probaron 6 intentos).

\section{Accesible para todos}
Se revisó la página con \textbf{axe-core}, la herramienta estándar para revisar accesibilidad (normas WCAG 2.1 AA), en cuatro
vistas: escritorio y celular, de día y de noche. La primera revisión encontró \textbf{67} fallas, la mayoría textos con poco
contraste. Se corrigieron hasta \textbf{0}. Además, la página se puede recorrer completa con el teclado.

\textbf{Lo que no se hizo}: el plan pedía probar la página con personas reales. Brandon eligió ``solo la revisión
automática'', y se declara que no hubo prueba con personas.
